% Rebuttal to Reviewer GBAF — CoLM 2026 Submission 991
% Compilable on Overleaf (pdflatex)
\documentclass[11pt]{article}
\usepackage[margin=1in]{geometry}
\usepackage{amsmath, amssymb}
\usepackage{enumitem}
\usepackage{parskip}
\usepackage{hyperref}
\hypersetup{colorlinks=true, linkcolor=blue, urlcolor=blue}

\title{Rebuttal to Reviewer GBAF\\
\large CoLM 2026, Submission 991}
\author{}
\date{}

\begin{document}
\maketitle

Dear Reviewer GBAF,

Thank you for the careful and supportive review. We are especially glad that the Section~5.2 finding, that strong correlation with human ratings does not by itself guarantee reliable point scoring, came through as one of the paper's central points, and that the choice of three heterogeneous judges (a specialist model, a larger reasoning model, and a closed-source generalist) read as the deliberate design it was rather than an arbitrary sample. Your two suggestions, a more formal task taxonomy and inline citations for the conformal methods, are both clear improvements, and we are happy to make them. We address each below and close with a short summary of the camera-ready revisions.

\section*{R1. A more formal task taxonomy}

You are right that the current four-way grouping (Vision-Heavy, General VQA, Knowledge-Web, Aesthetics-AI) is informal, and the COCO-versus-TextVQA placement you flagged is a clear example of the problem. We will replace it with a principled taxonomy drawn from existing work.

Concretely, we will re-categorize the 14 MLLM-as-a-Judge tasks using the six vision-capability dimensions from MM-Vet (Yu et al., 2023): Recognition, OCR, Knowledge, Language Generation, Spatial Awareness, and Math reasoning. MM-Vet is a widely cited multimodal benchmark whose taxonomy was built specifically to separate these capabilities. The placement you flagged resolves cleanly under it: COCO (natural-image captioning) maps to Recognition and Language Generation, while TextVQA (reading text from images) maps to OCR, so the two tasks are distinguished by the capability they actually exercise rather than by our informal labels. The apparent arbitrariness in calling one ``General VQA'' and the other ``Vision-Heavy'' disappears. The rest of the tasks map just as directly: MathVista to Math, ChartQA and InfographicsVQA to OCR and Math, ScienceQA to Knowledge.

For tasks that span more than one capability, we will assign a single primary capability, which is the one used for the per-task width row in Table~3 so that each task still appears exactly once with one width, and add a secondary-capabilities column for transparency. This keeps the table clean while recording the multi-capability structure. We will then rewrite \S 5.4 to describe the task-dependent uncertainty pattern under this grouping, reporting the resulting per-capability widths as they come out rather than predicting their shape in advance, and stating plainly how the task-dependent trend already visible in Table~3 appears once the rows are grouped by capability. MM-Vet (Yu et al., 2023) will be cited as the source. We note that only the grouping of rows changes: the per-task width values themselves, including the individual Table~3 rows and the per-dataset RSG analysis in \S 5.5, are unchanged.

\section*{Q1. Citations for the conformal methods, and tightening \S 5.1}

You are right that \S 5.1 lists the method names without inline citations, which assumes more familiarity with conformal prediction than we should. We will place a citation directly after each method name. The bibliography entries are:

\begin{itemize}[leftmargin=*]
\item Naive Split Conformal Prediction (Vovk et al., 2005)
\item Conformalized Quantile Regression, CQR (Romano et al., 2019)
\item Asymmetric CQR (Sesia \& Cand\`{e}s, 2020)
\item Conformal Histogram Regression, CHR (Sesia \& Romano, 2021)
\item Locally-adaptive Variance-based Distribution-free, LVD (Lin et al., 2021)
\item Boosted CQR and Boosted LCP (Xie et al., 2024)
\item Ordinal Adaptive Prediction Sets, OrdinalAPS (Lu et al., 2022)
\item Regression-as-Classification Conformal Prediction, R2CCP (Guha et al., 2024)
\end{itemize}

Most of these are already in the bibliography from citations elsewhere in the paper; we simply did not surface them inline in \S 5.1. This list also explains a count that could otherwise look inconsistent: Boosted CQR and Boosted LCP are two variants from a single source paper (Xie et al., 2024), so the paper speaks of eight conformal methods while Table~2 shows nine rows, the two Boosted variants being one family. We will make this explicit in the \S 5.1 prose so the count is unambiguous.

We also take seriously your point that this material contributes less to the main narrative and could be moved to the conclusion. We think the appendix is the more natural home for the full per-method detail, since it keeps the main text centered on the empirical findings while preserving the methodological content for readers who want it, but we are glad to use the conclusion instead if you would prefer. In practice the camera-ready will keep a short summary in \S 5.1, namely which methods we evaluate and why R2CCP and CHR are our headline choices, and move the full per-method descriptions and comparison details into the appendix.

\section*{Summary of camera-ready revisions}

In response to your review we will make the following changes:
\begin{enumerate}[leftmargin=*]
\item \textbf{\S 5.4 and Table~3.} Replace the informal four-way grouping with the MM-Vet six-capability taxonomy (Yu et al., 2023), re-aggregate the per-task uncertainty results under it using a primary capability per task with a secondary-capabilities column, and revise the \S 5.4 prose accordingly. The per-task width values are unchanged; only the grouping is re-labeled.
\item \textbf{\S 5.1.} Add an inline citation after every conformal-method name (Naive Split, CQR, Asymmetric CQR, CHR, LVD, Boosted CQR, Boosted LCP, OrdinalAPS, R2CCP), clarify the eight-methods-nine-rows count, compress the main-text prose to a short summary, and move the full per-method comparison into the appendix.
\end{enumerate}

We are grateful for the supportive review and for two suggestions that sharpen exactly the parts of the paper that needed it. We remain available throughout the discussion period for any further questions or to share intermediate materials if that would help.

\end{document}